\ifdefined\pdfsuppresswarningpagegroup%
  \pdfsuppresswarningpagegroup=1
\fi
\documentclass[11pt]{article}

\usepackage[a4paper,margin=25mm]{geometry}
\usepackage[T1]{fontenc}
\usepackage[utf8]{inputenc}
\usepackage{lmodern}
\usepackage{microtype}
\usepackage{amsmath,amssymb}
\usepackage{booktabs}
\usepackage{tabularx}
\usepackage{array}
\usepackage{graphicx}
\usepackage{caption}
\usepackage[numbers,sort&compress]{natbib}
\usepackage[hidelinks]{hyperref}
\usepackage{xcolor}

\setlength{\parindent}{0pt}
\setlength{\parskip}{0.55em}

\newcommand{\NJG}{\texttt{NJG\_MICHELL}}
\newcommand{\ProjectRecon}{\texttt{PROJECT\_RECONSTRUCTION}}
\newcommand{\DirectPlato}{\texttt{DIRECT\_PLATO\_TEXT}}
\newcommand{\MichellExplicit}{\texttt{MICHELL\_EXPLICIT}}
\newcommand{\ArchivePending}{\texttt{ARCHIVAL\_PENDING}}

% Publication figures are frozen first as white-background SVGs.
% The manuscript remains compilable before PDF derivatives exist.
\newcommand{\PaperFigure}[4]{%
  \begin{figure}[htbp]
    \centering
    \IfFileExists{#1}{%
      \includegraphics[width=#2\linewidth]{#1}%
    }{%
      \fbox{%
        \parbox[c][0.28\textheight][c]{0.90\linewidth}{%
          \centering
          \textbf{Publication figure placeholder}\\[0.7em]
          \texttt{\detokenize{#1}}\\[0.7em]
          The frozen white-background SVG exists in
          \texttt{paper/figures/}; a PDF derivative will be inserted at
          manuscript-assembly time.
        }%
      }%
    }
    \caption{#3}%
    \label{#4}
  \end{figure}
}

\title{\textbf{Reconstructing John Michell's New Jerusalem Geometry}\\
\large A Reproducible Audit of Euclidean Structure, Sevenfold Constructibility,
and the Plato\textendash{} Michell Whorl Interpretation}

\author{
Salah-Eddin Gherbi\\
Independent Researcher, United Kingdom\\
ORCID\@: \mbox{0009\mbox{-}0005\mbox{-}4017\mbox{-}1095}
}

\date{}

\begin{document}

\maketitle

\begin{abstract}
We present a reproducible computational and historical audit of John Michell's
New Jerusalem geometry.  The analysis separates source statements, project
reconstructions, exact mathematical results, approximate historical
constructions, source-calibrated comparisons, and downstream visualisation.
A unit-only generative model reconstructs the principal Earth\textendash{} Moon,
twelve-circle, outer-wall, sevenfold, and 28-point components without feeding
measured source-plate coordinates back into the geometric builder.  Direct
calibration of Michell's printed figures supports a near-regular
$\{7/2\}$ topology for Figure~14 and a polar-pivot project reconstruction for
the Figure~12 outer wall, while preserving the distinction between
source-consistency and independent prediction.  Michell's two approximate
sevenfold procedures are reconstructed separately; their local steps bracket
the exact $2\pi/7$ angle, and the second method has the exact project-derived
form $\alpha_2=\arccos(5/8)$.  Algebraically,
$y=2\cos(2\pi/7)$ satisfies the irreducible cubic
$y^3+y^2-2y-1=0$, excluding exact regular sevenfold construction from the
frozen ordinary straightedge-and-compass closure.  Within the registered
operation-count model, zero cubic-capable operations are impossible and one
angle trisection is sufficient.  Finally, a source audit of Plato's
\emph{Republic} and Michell's later numerical interpretation of the eight
whorls reproduces Michell's printed arithmetic with zero fitted parameters,
without attributing his numerical widths to Plato or inferring historical
transmission from antiquity.
\end{abstract}

\section{Introduction}

John Michell's New Jerusalem diagrams combine a scale-free geometric core with
historical metrology, approximate polygonal constructions, and later numerical
interpretation \citep{michell1973city,michell2009world}.  Their visual
complexity makes several logically different questions easy to conflate:
what Michell explicitly states, what a printed plate actually depicts, what can
be reconstructed from explicit Euclidean rules, what is mathematically exact,
and what remains interpretive.

The present work treats those questions as separate evidential layers.
Rather than fitting a single picture until it resembles a source image, the
project freezes construction rules, source witnesses, calibration protocols,
and downstream comparisons at explicit checkpoints.  The resulting
\NJG{} composite is generated from the model scale alone; source digitisation
is used for evaluation rather than as a hidden geometric input.

The paper has four principal aims.  First, it records the deterministic
New Jerusalem reconstruction and the source-calibration boundary for
Figures~12 and~14.  Second, it compares Michell's two approximate sevenfold
methods without combining or optimising them.  Third, it identifies the exact
algebraic constructibility boundary of the sevenfold layer and a minimal
one-cubic-operation extension within the frozen project grammar.  Fourth, it
separates Plato's ordinal description of the eight whorls from Michell's
numerical assignment and reproduces Michell's arithmetic without back-projecting
those numbers into Plato.

\section{Sources and methods}

\subsection{Evidence hierarchy}

The analysis uses explicit provenance classes rather than a single undifferentiated
notion of ``agreement''.  Table~\ref{tab:provenance} summarises the principal
classes used in the manuscript.

\begin{table}[ht]
\centering
\caption{Evidence and provenance classes used throughout the audit.}%
\label{tab:provenance}
\begin{tabularx}{\linewidth}{>{\ttfamily}l X}
\toprule
Class & Meaning \\
\midrule
DIRECT\_PLATO\_TEXT &
Literal information recoverable from the frozen Plato witnesses. \\
MICHELL\_EXPLICIT &
Numerical or construction statements explicitly supplied by Michell. \\
PROJECT\_RECONSTRUCTION &
Deterministic geometry or arithmetic carried out by the project. \\
source-calibrated comparison &
Comparison against digitised printed source geometry. \\
exact mathematical result &
Algebraic or Euclidean result established independently of visual fit. \\
ARCHIVAL\_PENDING &
Historical lead not established from a frozen primary witness. \\
\bottomrule
\end{tabularx}
\end{table}

\subsection{Historical witnesses}

The core Michell source plate used for Figures~12 and~14 is the 1973 Abacus
edition of \emph{City of Revelation} \citep{michell1973city}.  Figure~194 and
the later Sommerville-attributed dimensional discussion are audited from
\emph{How the World Is Made} \citep{michell2009world}.  Michell's numerical
interpretation of Plato's eight whorls is taken from the purchased 2008 digital
edition of \emph{The Dimensions of Paradise} \citep{michell2008dimensions}.

For Plato, the project freezes both English and Greek Perseus TEI witnesses
for \emph{Republic}~616d--617b.  The underlying English translation is Shorey's
Loeb edition \citep{plato1935republic}; the Greek text derives from Burnet's
Oxford edition \citep{plato1902republic}.  The exact project XML files are
pinned to one upstream Perseus commit and hashed independently
\citep{perseusRepublicWitnesses}.  The digital-witness identity is therefore
kept separate from the bibliographic identity of the underlying print
editions.

The external 2008 Wikimedia diagram is included only as a reproducible
historical computational comparator, \texttt{NJG\_SVG}; it is not treated as
a Michell primary source \citep{anonmoos2008newjerusalem}.

\subsection{Normalized geometry and deterministic generation}

The normalized cardinal core uses Earth radius $5.5$, Moon radius $1.5$,
construction-circle radius $7$, and Earth-square side $11$.  The complete
\NJG{} model adds the twelve-Moon system, the preferred outer-wall project
reconstruction, Michell's approximate Method-1 sevenfold construction,
the reciprocal fourteenfold construction, the four-triangle 28-point
scaffold, and the scaffold-derived Figure~14 candidate.

Only the overall scale unit is caller supplied.  The source-calibration
modules remain downstream of the generative builder, so measured plate
coordinates do not refit the model.

\subsection{Source-plate calibration}

Printed source geometry is treated as measured historical evidence rather
than an exact Euclidean object.  Repeated landmark digitisation is followed
by registration-model comparison and frozen no-refit candidate evaluation.
For Figure~14, edge topology is tested independently from endpoint location
because regular $\{7/2\}$ and $\{7/3\}$ heptagrams share the same seven-vertex
set.

\subsection{Exact constructibility analysis}

The exact sevenfold target is represented algebraically by
\begin{equation}
y=2\cos\left(\frac{2\pi}{7}\right),
\end{equation}
whose minimal polynomial is
\begin{equation}
y^3+y^2-2y-1=0.
\end{equation}
The degree-three obstruction is compared with the quadratic-extension closure
of ordinary straightedge-and-compass constructions.  A separate,
preregistered project extension then permits exactly one cubic-capable
\texttt{TRISECT\_ANGLE} operation.

\section{Core New Jerusalem reconstruction}

\PaperFigure{figures/pdf/figure01_njg_michell_composite.pdf}
  {0.76}
  {Integrated unit-only \NJG{} reconstruction.  The publication derivative is
  downstream of the frozen generative composite; source calibration does not
  feed geometric coordinates back into the builder.}
  {fig:njg-composite}

Figure~\ref{fig:njg-composite} summarises the integrated construction.
The distinction between the generative object and the later evidence layers is
central: source consistency can evaluate a frozen candidate without becoming
a hidden fitting channel.

\subsection{Figure 12 outer wall}

Three wall rules were compared against the affine-calibrated Figure~12 plate.
The preferred \texttt{POLAR\_PIVOT\_TANGENT} reconstruction gives the best
wall-direction agreement among the registered candidates, but the
support-distance metric alone favours another comparator.  It is therefore
described as a source-supported project reconstruction rather than a uniquely
established historical construction.

\begin{table}[ht]
\centering
\caption{Selected frozen Figure~12 diagnostics for the preferred polar-pivot
wall reconstruction.}%
\label{tab:figure12}
\begin{tabular}{lr}
\toprule
Metric & Frozen value \\
\midrule
Wall-normal angular RMS & $1.019836^\circ$ \\
Wall-support RMS & $0.038999\,u$ \\
Wall-vertex RMS & $0.100022\,u$ \\
Perimeter difference & $-0.379134\%$ \\
Area difference & $-0.709332\%$ \\
\bottomrule
\end{tabular}
\end{table}

\clearpage

\section{Source-calibrated sevenfold geometry}

\subsection{Figure 14 topology}

Repeated endpoint and stroke tracing supports a near-regular
$\{7/2\}$ heptagram topology in Michell's Figure~14
\citep{michell1973city}.  This supersedes the project's provisional
$\{7/3\}$ interpretation.  Figure~\ref{fig:figure14} presents the frozen
source-calibrated comparison.

\PaperFigure{figures/pdf/figure02_figure14_heptagram.pdf}
  {0.95}
  {Source-calibrated Figure~14 comparison.  Endpoint regularity and printed
  stroke topology are treated as separate evidence channels; the calibrated
  topology supports $\{7/2\}$.}
  {fig:figure14}

\subsection{Michell Method 1: \texorpdfstring{\(7\rightarrow14\rightarrow28\)}{7 to 14 to 28}}

Michell's first Figure~194 branch generates an approximate sevenfold step,
its reciprocal fourteenfold construction, and the four-triangle 28-point
scaffold \citep{michell2009world}.  Figure~\ref{fig:method1} shows the frozen
Method~1 scaffold used by the integrated construction.

\PaperFigure{figures/pdf/figure03_method1_28_point_scaffold.pdf}
  {0.95}
  {Michell Method~1 and the frozen 28-point scaffold used in the integrated
  construction.}
  {fig:method1}

\subsection{Michell Method 2: \texorpdfstring{\(7\rightarrow21\rightarrow42\)}{7 to 21 to 42}}

The second Figure~194 branch is kept distinct.  From the source-described
equilateral-triangle and side-bisecting-arc construction, the project derives
\begin{equation}
\alpha_2=\arccos\left(\frac{5}{8}\right).
\end{equation}
The complete 21- and 42-point nonuniformity follows from the local deficit
\begin{equation}
\delta=\frac{2\pi}{7}-\alpha_2.
\end{equation}
The symbolic gap identities are post-result explanatory derivations rather
than preregistered confirmations.  Figure~\ref{fig:method2} summarizes the
registered Method~2 local construction and its deterministic 21- and
42-point propagation.

\PaperFigure{figures/pdf/figure04_method2_7_21_42.pdf}
  {0.98}
  {Michell Method~2: local seven-mark construction and deterministic
  propagation to 21 and 42 points.}
  {fig:method2}

\begin{table}[ht]
\centering
\caption{Exact sevenfold step and Michell's two frozen approximations.}%
\label{tab:sevenfold-comparison}
\begin{tabular}{lrrr}
\toprule
Construction & Step (deg) & Signed error (deg) & Seven-step closure (deg) \\
\midrule
Exact & 51.428571428571 & 0 & 0 \\
Method 1 & 51.470701432440 & +0.042130003869 & +0.294910027080 \\
Method 2 & 51.317812546511 & -0.110758882061 & -0.775312174426 \\
\bottomrule
\end{tabular}
\end{table}

Thus the two source-described approximations bracket the exact value:
\[
\text{Method 2}<\text{exact}<\text{Method 1}.
\]

\section{Exact sevenfold boundary}

Because the target coordinate has algebraic degree three over
$\mathbb{Q}$, the exact regular sevenfold is excluded from the frozen ordinary
straightedge-and-compass closure.  This conclusion is theorem-level and does
not depend on failure to find a sufficiently accurate numerical candidate.

After that boundary was frozen, the project registered one cubic-capable
operation.  With the native dimensions $R=7$ and $D=3$,
\[
u=R-2D=1,\qquad
v=D\sqrt{3}=3\sqrt{3},\qquad
H=2\sqrt{7}.
\]
For the resulting anchor angle $\Theta$,
\[
\cos\Theta=\frac{1}{2\sqrt{7}}.
\]
One angle trisection sets $\phi=\Theta/3$.  Defining
\[
z=\frac{2\sqrt{7}}{3}\cos\phi
\]
and applying the triple-angle identity gives
\[
z^3-\frac73z-\frac7{27}=0.
\]
With $y=z-\frac13$, this becomes exactly
\[
y^3+y^2-2y-1=0.
\]

Within the frozen operation-count model, the combined result is therefore
\[
0\ \text{cubic-capable operations: impossible},\qquad
1\ \text{cubic-capable operation: sufficient}.
\]
This result is a project construction and is not attributed to Michell,
Sommerville, Plato, or an ancient source.  Figure~\ref{fig:exact-heptagon}
places the exact construction beside the frozen Method~1 and Figure~14
comparators without rotational refitting.

\PaperFigure{figures/pdf/figure05_exact_heptagon_synthesis.pdf}
  {0.98}
  {Exact-heptagon comparison and one-trisection synthesis.  The exact
  construction is shown alongside the frozen Method~1 and Figure~14
  comparators without rotational fitting or endpoint reassignment.}
  {fig:exact-heptagon}

\clearpage

\section{Historical metrology and Sommerville-attributed comparisons}

Historical dimensional agreement is classified according to whether a source
quantity was a construction input, a development-used target, a retrospective
comparison, a derived relation, or an unevaluated archival claim.  Development
targets are not later relabelled as independent predictions.

Michell reports dimensional results attributed to Ian Sommerville in
\emph{How the World Is Made} \citep{michell2009world}.  The underlying 1974
primary item has not been frozen in the present project and remains
\ArchivePending{}.

\begin{table}[ht]
\centering
\caption{Frozen no-refit dodecagon dimensional comparisons.  Derived and
equivalent rows are not additional independent observations.}%
\label{tab:sommerville}
\begin{tabular}{lrrr}
\toprule
Quantity & Historical value & Frozen value & Residual \\
\midrule
Long radius & 6336 ft & 6335.8902241097 ft & $-0.001732574\%$ \\
Short radius & 6300 ft & 6302.1210766004 ft & $+0.033667883\%$ \\
Class ratio & $176/175$ & 1.0053583780919 & $-0.035388542\%$ \\
\bottomrule
\end{tabular}
\end{table}

These values are retrospective source-consistency checks.  They do not
establish independent verification of the unresolved 1974 primary paper,
statistical significance, ancient intentionality, or physical law.

\section{Plato's eight whorls and Michell's numerical interpretation}

\subsection{Literal Plato layer}

The frozen \emph{Republic} passage describes eight nested whorls and gives
their rim widths in the ordinal order
\[
1>6>4>8>7>5>3>2
\]
\citep{plato1935republic,plato1902republic,perseusRepublicWitnesses}.
The audited passage supplies no numerical magnitudes or ratios for those rim
widths.  Its speed ordering is a separate textual attribute and is not used as
a substitute for the rim-width order.

\subsection{Michell's numerical layer}

Michell preserves the same ordinal width ordering but supplies the musical
number set
\[
6,\ 8,\ 9,\ 12,\ 18,\ 24,\ 27,\ 36
\]
together with the shaft rule, the centre-out ring ordering, and a scale factor
of $180$ \citep{michell2008dimensions}.  Assigning the distinct numerical
values according to the frozen ordinal ranking produces the centre-out widths
\[
18,\ 12,\ 27,\ 9,\ 24,\ 8,\ 6,\ 36.
\]

The shaft rule gives
\[
6\times\frac23=4,
\]
and successive addition gives cumulative radii
\[
4,\ 22,\ 34,\ 61,\ 70,\ 94,\ 102,\ 108,\ 144.
\]
Scaling by $180$ yields
\[
720,\ 3960,\ 6120,\ 10980,\ 12600,\ 16920,\ 18360,\ 19440,\ 25920.
\]

\begin{table}[ht]
\centering
\caption{Frozen centre-out Plato\textendash{} Michell whorl reconstruction.}%
\label{tab:whorls}
\begin{tabular}{rrrr}
\toprule
Centre-out whorl & Width & Cumulative radius & Scaled radius \\
\midrule
shaft & \textemdash{} & 4 & 720 \\
8 & 18 & 22 & 3960 \\
7 & 12 & 34 & 6120 \\
6 & 27 & 61 & 10980 \\
5 & 9  & 70 & 12600 \\
4 & 24 & 94 & 16920 \\
3 & 8  & 102 & 18360 \\
2 & 6  & 108 & 19440 \\
1 & 36 & 144 & 25920 \\
\bottomrule
\end{tabular}
\end{table}

The reconstruction uses zero free parameters, zero permutation searches, zero
scale fits, zero optimisations, and zero nearest-match choices.  Exact
reproduction validates the internal numerical chain of Michell's printed
construction.  It does not establish that Plato supplied those numerical
widths, intended Michell's assignment, or transmitted the scheme historically
from antiquity.  Figure~\ref{fig:whorls} summarizes the radial construction
while keeping the Plato, Michell, and project provenance layers explicit.

\PaperFigure{figures/pdf/figure06_plato_michell_whorls.pdf}
  {0.98}
  {The Spindle of Necessity---Plato's Eight Whorls.  The figure separates
  Plato's ordinal information, Michell's numerical interpretation, and the
  project's deterministic arithmetic reconstruction.}
  {fig:whorls}

\clearpage

\section{Discussion}

The audit produces three different kinds of positive result that should not be
collapsed into one another.

First, deterministic reconstruction establishes that a large part of the New
Jerusalem diagram can be generated from explicit geometric rules with only an
overall scale parameter.  Second, source calibration establishes which frozen
project candidates are more consistent with particular printed plates, but
does not transform those source-dependent comparisons into statistically
independent predictions.  Third, the exact sevenfold result is mathematical:
the degree-three obstruction and the one-trisection construction do not depend
on the quality of the historical plate reproductions.

The historical analysis is similarly bounded.  Plato supplies an ordinal
eight-whorl structure; Michell supplies a later numerical interpretation; and
the project can reconstruct the latter exactly.  Numerical reproducibility is
not evidence for ancient transmission.

Several limitations remain.  Printed-source calibration is subject to
reproduction quality, line thickness, manual digitisation, and registration
uncertainty.  Historical metrology includes quantities that entered model
development and therefore cannot be recycled as independent validation.
The Sommerville 1974 primary source remains archival-pending.  Finally, the
one-trisection result establishes minimality only in the number of
cubic-capable operations within the frozen project grammar; it does not show
that trisection is uniquely simplest among all non-Euclidean construction
technologies.

\section{Conclusion}

A provenance-aware computational treatment makes it possible to separate what
Michell states, what his printed figures support, what the project reconstructs,
and what exact mathematics establishes.  The resulting New Jerusalem model is
reproducible and source-auditable while retaining explicit boundaries around
historical interpretation.  The same framework resolves the topology of the
Figure~14 heptagram, characterises Michell's two approximate sevenfold methods,
identifies the exact constructibility boundary and a one-trisection extension,
and reconstructs Michell's Plato-whorl arithmetic without attributing his
numerical scheme to Plato.

\section*{Data and code availability}

The complete computational record, frozen numerical outputs, source-provenance
manifests, deterministic figure generators, and automated tests are maintained
in the \emph{New Jerusalem Geometry} repository.

\textbf{Release-boundary note.}
The manuscript is being prepared while the package metadata deliberately
remains at version \texttt{0.9.0}.  The final public repository reference,
version \texttt{1.0.0}, archival Zenodo DOI, and corresponding software
bibliography entry will be inserted only after the controlled public release.

Copyrighted Michell source scans and purchased ebook screenshots are not
redistributed.  Their identities are preserved by local source manifests and
cryptographic hashes.  The frozen Perseus XML witnesses retain their upstream
third-party licensing boundary.

\begingroup
\raggedright{}
\bibliographystyle{plainnat}
\bibliography{references}
\endgroup

\end{document}
